\documentclass[a4paper,12pt]{jsarticle}
\usepackage{geometry}
\geometry{top=25mm, bottom=25mm, left=25mm, right=25mm}
\usepackage{booktabs}
\usepackage{listings}
\usepackage{xcolor}
\usepackage{hyperref}
\hypersetup{colorlinks=true, linkcolor=black, urlcolor=blue}

% ---- listings 設定（Processing は Java ベース）----
\lstdefinestyle{processing}{
  language=Java,
  basicstyle=\ttfamily\small,
  keywordstyle=\color{blue}\bfseries,
  commentstyle=\color{gray}\itshape,
  stringstyle=\color{red!60!black},
  numbers=left,
  numberstyle=\tiny\color{gray},
  breaklines=true,
  frame=single,
  tabsize=2,
  captionpos=b,
  showstringspaces=false
}

\title{チーム23 ハッカソン1\\
       基本設計・詳細設計\\[0.5em] 
       }     
\author{梅澤颯太（学籍番号：25G1021）}
\date{\today}

\begin{document}
\maketitle
\tableofcontents
\newpage

\section{基本設計}

\subsection{117：音色合成 基本設計}

\subsubsection{目的}
Processingを使って，金管楽器らしい音色を
作る方針を決める．
「倍音の重ね合わせ（加算合成）」と「振幅エンベロープ（音の大きさの時間変化）」
を組み合わせる．

\subsubsection{音色の仕組み}
楽器音の音色は以下の2つで決まる．

\begin{itemize}
  \item \textbf{倍音の含み方}：
        基音（出したい音の周波数）の整数倍の周波数成分をいくつ，どの強さで重ねるか．
        偶数倍音・奇数倍音の有無や強さが音色の違いに大きく影響する．
  \item \textbf{倍音の立ち上がり方（エンベロープ）}：
        音が鳴り始めてから消えるまでの音量の変化．
        楽器らしさはここで決まることが多い．
\end{itemize}

\noindent
\texttt{WavetableGenerator.gen10()} を使って倍音の振幅を指定する．

\subsubsection{倍音構成の方針}

金管楽器の特徴は「倍音が多く，明るくて力強い音」である．
具体的には，基音（第1倍音）から第2・第3倍音あたりが特に強く，
高次の倍音になるほど徐々に弱くなる．

\begin{table}[h]
  \centering
  \caption{金管楽器風の倍音構成（案）}
  \label{tab:harmonics}
  \begin{tabular}{ccl}
    \toprule
    倍音次数 & 振幅の大きさ（0.0〜1.0） & 根拠・メモ \\
    \midrule
    第1倍音（基音） & 1.00 & 一番大きい成分 \\
    第2倍音         & 0.70 & 金管では2倍音が強い \\
    第3倍音         & 0.50 & 明るさを加える \\
    第4倍音         & 0.30 & 厚みを加える \\
    第5倍音         & 0.20 & 少しブライトにする \\
    第6倍音         & 0.10 & 高次はだんだん弱く \\
    第7倍音         & 0.05 & ごくわずかに含む \\
    \bottomrule
  \end{tabular}
\end{table}

\subsubsection{エンベロープの方針}

授業で使った \texttt{HackInstrument} クラスでは，\texttt{Line} を使って
「音が鳴り始めから時間をかけて音量をゼロに下げていく（減衰）」という処理をしている．
金管楽器らしくするためには，さらに「鳴り始め（立ち上がり）」も加えたい．

\begin{description}
  \item[立ち上がり（Attack）] 音が鳴り始めてから最大音量に達するまでの時間．
        約 20 ミリ秒（とても短い）．
  \item[減衰・持続（Decay / Sustain）] 最大音量から少し落ち着いた音量で安定する．
        息を吹き続けている感じ．
  \item[余韻（Release）] 息を止めて音が消えるまでの時間．約 60 ミリ秒．
\end{description}

\noindent
ただし，Minim の \texttt{Line} で表現できるのは「1本の直線的な変化」のみである．
そのため音の出来次第でADSRクラスを使う方針とする．

\section{詳細設計}

\subsection{124：Arduino からの NOTE 受信と発音の仕組み}
\label{sec:detail-124}

\subsubsection{概要}
Arduinoから Serial 通信で届く音の指令を
Processing が受け取り，対応する音を鳴らす．
通信フォーマットの詳細は115と130が
確定してから決める．ここでは大まかな流れのみ示す．

\subsubsection{処理の流れ（案）}

\begin{enumerate}
  \item Arduino が Serial 通信で「どの音を・どの長さで・どの大きさで鳴らすか」
        という情報を Processing に送る．
  \item Processing の \texttt{serialEvent()} 関数が自動で呼ばれ，データを受け取る．
  \item 受け取ったデータから音名・長さ・振幅を取り出す．
  \item 取り出した情報をもとに \texttt{HackInstrument} を生成し，
        \texttt{out.playNote()} で音を鳴らす．
\end{enumerate}

\subsection{125：4パート発音管理}
\label{sec:detail-125}

4つの楽器パート（現段階では金管楽器）を
それぞれ独立して鳴らせるように管理する仕組みを作る．
授業では，melody・duration・startTime のような配列を使って
音の情報を管理する方法を学んだ．
同じ発想で，パートごとの音色（\texttt{Waveform}）を配列で管理する．

\subsection{126：金管楽器風音色の具体実装}
\label{sec:detail-126}

\subsubsection{概要}
タスク117で決めた倍音構成とエンベロープを，
Processing のコードに実際に落とし込む．

\subsubsection{倍音テーブルの実装（案）}

授業で使った \texttt{WavetableGenerator.gen10()} にそのまま値を入れる．

\begin{lstlisting}[style=processing, caption={金管楽器風の波形生成（案）}]
Waveform makeBrassWaveform() {
  return WavetableGenerator.gen10(
    4096,  // サンプル数（2の累乗にすること）
    new float[] {
      1.00f,  // 第1倍音（基音）：最大
      0.70f,  // 第2倍音：金管は2倍音が強い
      0.50f,  // 第3倍音
      0.30f,  // 第4倍音
      0.20f,  // 第5倍音
      0.10f,  // 第6倍音
      0.05f   // 第7倍音：ほんの少し
    }
  );
}
\end{lstlisting}

\subsubsection{エンベロープの実装（案）}

既存の \texttt{HackInstrument} クラスの \texttt{noteOn()} 内で
\texttt{ampEnv.activate(duration, maxAmp, 0)} を呼ぶ．
これにより，音符の長さ全体をかけて最大音量からゼロへ直線的に減衰する．鳴り始めが最大音量で，そこから徐々に消えていく形になる．
\texttt{Line} は一度に一本の直線変化しか持てないため，
立ち上がり（Attack）の表現は行わない．（音の出来次第でADSRクラスを使う）

 
\begin{table}[h]
  \centering
  \caption{\texttt{ampEnv.activate()} に渡す値}
  \label{tab:line-params}
  \begin{tabular}{ccc}
    \toprule
    時間 & 開始音量 & 終了音量 \\
    \midrule
    \texttt{duration}（音符の長さ全体） & \texttt{maxAmp} & 0.0 \\
    \bottomrule
  \end{tabular}
\end{table}

 
\subsubsection{パラメータ調整の方針}
 
表\ref{tab:harmonics}・表\ref{tab:line-params}の値はいずれも仮置きである．
実際に音を鳴らしながら耳で確認し，金管楽器らしく聞こえるよう調整する．
127（FFT解析）の結果が得られた場合はその値に更新する．

\subsection{127：Python による FFT 解析}
\label{sec:detail-127}

\subsubsection{概要}
実際の金管楽器の音声ファイルを Python で周波数分析（FFT）し，
倍音の強さとエンベロープの時定数を数値として取り出す．
取り出した値を表\ref{tab:harmonics}および \texttt{attackTime} に反映する．

\subsubsection{手順（案）}

\begin{enumerate}
  \item フリー素材から金管楽器の単音サンプルを入手する．
  \item Python の \texttt{scipy.fft} または \texttt{librosa} で FFT を実施し，
        各倍音成分のパワーを求める．
  \item 基音のパワーを 1.0 として正規化し，第2〜第7倍音の比率を計算する．
  \item 音量の時間変化（RMS カーブ）から立ち上がり時間を推定する．
  \item 得られた値を \texttt{makeBrassWaveform()} の配列と
        \texttt{attackTime} に反映する．
\end{enumerate}

\section{生成AIの利用について}
本レポートを作成するにあたり，ClaudeAIを利用しました．生成された内容を吟味しつつ，修正・加筆して完成させました．

\end{document}